%
% Chapter 6: Conclusion
%

\chapter{Conclusion and Future Directions}
\label{chap:conclusion}

\section{Summary of Contributions}

This dissertation connects an archival automated preparation pilot to a deployable \textbf{Flowmatic} platform:
\begin{itemize}
  \item six-dimensional DQI methodology and classical preparation \emph{design} (June 2025 offline pilot; not re-run in-repo);
  \item Dockerised batch and smart-city pipelines with reproducible API evidence;
  \item an eight-model neural portfolio with multi-seed held-out test metrics on Hugging Face;
  \item adaptive Manual/Auto core-unit routing over a documented model registry;
  \item open Phase~2--3 scripts (\texttt{run\_phase2\_prep.py}, \texttt{run\_phase3\_core.py}) without embedding classical XGBoost in production inference.
\end{itemize}

\section{Key Findings}

Platform preparation on 30\,000 Astana upload records completes in 406\,ms with full quality scoring. The neural severity classifier reaches macro-F1 = $0.911 \pm 0.017$ under a shared temporal hold-out protocol. Archival tabular preparation gains (XGBoost pilot) motivated the workflow but are not merged into the HF checkpoint line.

\section{Limitations}

\begin{itemize}
  \item Astana data remain semi-synthetic; live agency validation is limited to the workbench demo.
  \item External graph benchmarks (METR-LA, PEMS-BAY) are ingested but not yet evaluated as hold-out targets for Astana-trained models.
  \item Router oracle ablations, CRPS, and AUROC event metrics are not yet reported.
  \item LLM copilot summaries are optional and not formally calibrated; SHAP applies only to the archival classical pilot.
  \item iTransformer speed forecasting under the current training budget is not production-ready (RMSE $\approx 1.0$).
\end{itemize}

\section{Future Work}

Planned extensions include Astana$\rightarrow$PeMS transfer experiments, auto-router confusion analysis, probabilistic forecasting metrics, operator-facing calibration studies, and retraining or retiring underperforming registry slots. Federated coordination remains a demonstration path pending production hardening.

\section{Concluding Remarks}

Automated preparation remains the foundation of trustworthy ITS analytics. Flowmatic shows that preparation stages can be implemented as an adaptive, multimodal platform whose \textbf{neural and operational} results are reproducible from open scripts and published checkpoints, while classical tabular experiments remain an explicitly scoped archival line documented in Table~\ref{tab:gap-analysis}.
